\section{Jordan 分割定理}\label{ch10:section-jordan-separation}

关于平面拓扑学, 有几个较为困难的问题, 它们在分析的研究中是很自然就产生的. 这些问题的答案看上去很直观, 但要给出证明却非常困难. 这些问题包括 Jordan 曲线定理, Brouwer 区域不变性定理, 以及简单闭曲线的环绕数 (winding number) 是 $0$ 或 $\pm 1$ 这一经典定理. 本章中, 我们运用覆叠空间和基本群理论来证明这些结论.

我们先考虑这些经典数学定理中的 Jordan 曲线定理. 这个定理涉及这样一个在几何上看似很直观的结论, 即平面上的一条简单闭曲线总会将平面分成两个连通分支, 一个在``内'', 一个在``外''. 这个猜测最初由 Camille Jordan 于 1892 年提出, 随后便出现了一些不正确的证明, 其中包括 Jordan 本人给出的一个证明. 正确的证明最终由 Oswald Veblen 于 1905 年给出. 早期的证明是很复杂的, 数年之后才找到了较为简单的证明方法. 如果借助于现代代数拓扑, 尤其是奇异同调论, 证明会变得很直接. 这里给出的证明是我们所知道的证明中最简单的一个, 只用到了覆叠空间和基本群理论中的结论.

我们将 Jordan 曲线定理的证明分为三部分. 第一部分称为 Jordan 分割定理: 用平面上的一条简单闭曲线来分割平面, 这时平面至少会被分成两个连通分支. 第二部分是说平面上的一条弧不能将平面分开. 第三部分为严格的 Jordan 曲线定理: 平面上的一条简单闭曲线 $C$ 将平面分为两个连通分支, 并且 $C$ 是这两个分支的共同边界. 本节就来讨论定理证明的第一部分.

在处理分割定理的过程中, 将所涉及的诸集合视为 $S^{2}$ 的子集较之视为 $\mathbb{R}^{2}$ 的子集会更为方便. 在证明了关于 $S^{2}$ 的分割定理之后, 关于 $\mathbb{R}^{2}$ 的分割定理便将顺理成章地得到. 两个定理中所涉及的集合之间的关系将由下述引理给出.

以前提过, 若 $b$ 是 $S^{2}$ 的一个点, 则存在着一个同胚 $h\colon S^{2}-b\to\mathbb{R}^{2}$. 例如, 可以先选取 $S^{2}$ 上的旋转, 将 $b$ 映射到北极, 接着再做球极投影.

\begin{lemma}[component lemma]\label{ch10:lemma-component}
设 $C$ 是 $S^{2}$ 的一个紧致子空间, $b$ 是 $S^{2}-C$ 的一个点, $h\colon S^{2}-b\to\mathbb{R}^{2}$ 是一个同胚. 又设 $U$ 是 $S^{2}-C$ 的一个分支. 如果 $U$ 不包含 $b$, 那么 $h(U)$ 是 $\mathbb{R}^{2}-h(C)$ 的一个有界分支. 如果 $U$ 包含 $b$, 那么 $h(U-b)$ 是 $\mathbb{R}^{2}-h(C)$ 的一个无界分支.

特别地, 如果 $S^{2}-C$ 有 $n$ 个分支, 那么 $\mathbb{R}^{2}-h(C)$ 便有 $n$ 个分支.
\end{lemma}

\begin{proof}
我们先来证明当 $U$ 是 $S^{2}-C$ 的一个分支时 $U-b$ 是连通的. 当 $b\notin U$ 时, 这个结论是平凡的. 因此我们假设 $b\in U$, 并且集合 $A,B$ 构成了 $U-b$ 的一个分割. 选取 $b$ 的一个邻域 $W$ 使得 $W$ 与集合 $C$ 无交, 并且 $W$ 同胚于 $\mathbb{R}^{2}$ 中的一个开球. 因为 $W$ 是连通的, 所以它包含于 $U$ 中. 因为 $W-b$ 是连通的, 所以它或者包含于 $A$, 或者包含于 $B$. 设 $W-b\subset A$. 由于 $W$ 是 $b$ 的一个与 $B$ 无交的邻域, 所以 $b$ 不是 $B$ 的一个极限点. 从而集合 $A\cup\{b\}$ 与 $B$ 构成了 $U$ 的一个分割, 与假设矛盾.

设 $\{U_{\alpha}\}$ 是 $S^{2}-C$ 的所有分支构成的集合. 令 $V_{\alpha}=h(U_{\alpha}-b)$. 由于 $S^{2}-C$ 是局部连通的, 所以 $U_{\alpha}$ 都是 $S^{2}$ 中的连通开子集, 并且两两无交. 从而 $V_{\alpha}$ 都是 $\mathbb{R}^{2}-h(C)$ 中的连通开子集, 并且两两无交, 因此 $V_{\alpha}$ 是 $\mathbb{R}^{2}-h(C)$ 的分支.

现将同胚 $h\colon S^{2}-b\to\mathbb{R}^{2}$ 扩充为同胚 $H\colon S^{2}\to\mathbb{R}^{2}\cup\{\infty\}$, 其中 $\mathbb{R}^{2}\cup\{\infty\}$ 是 $\mathbb{R}^{2}$ 的单点紧致化, 并令 $H(b)=\infty$. 如果 $U_{\beta}$ 是 $S^{2}-C$ 中包含 $b$ 的一个分支, 那么 $H(U_{\beta})$ 是 $\mathbb{R}^{2}\cup\{\infty\}$ 中包含 $\infty$ 的一个邻域. 从而 $V_{\beta}$ 是无界的. 因为它的补 $\mathbb{R}^{2}-V_{\beta}$ 是紧致的, $\mathbb{R}^{2}-h(C)$ 的所有其他分支都是有界的. 参见 \cref{ch10:fig-61-1}.
\end{proof}

\par
\begin{figure}[htbp]
\centering
\includegraphics[width=0.7\textwidth]{ch10/fig-61-1.png}
\caption{球极投影保持分支数.}\label{ch10:fig-61-1}
\end{figure}

\begin{lemma}[nullhomotopy lemma]\label{ch10:lemma-nullhomotopy}
设 $a$ 和 $b$ 是 $S^{2}$ 中的两个点, $A$ 是一个紧致空间,
\[
f\colon A\to S^{2}-a-b
\]
是一个连续映射. 如果 $a$ 和 $b$ 属于 $S^{2}-f(A)$ 的同一分支, 那么 $f$ 是零伦的.
\end{lemma}

\begin{proof}
我们用 $\mathbb{R}^{2}$ 的单点紧致化 $\mathbb{R}^{2}\cup\{\infty\}$ 代替 $S^{2}$, 并且分别将点 $a$ 和 $b$ 取作 $0$ 和 $\infty$. 这时引理化做: 设 $A$ 是一个紧致空间, $g\colon A\to\mathbb{R}^{2}-0$ 是一个连续映射. 若 $0$ 属于 $\mathbb{R}^{2}-g(A)$ 的无界分支, 则 $g$ 是零伦的.

上述结论易于证明. 选取球心在原点、半径充分大的球 $B$, 使得 $B$ 包含 $g(A)$. 设 $p$ 是 $\mathbb{R}^{2}$ 上不属于 $B$ 的一个点, 那么 $0$ 和 $p$ 都属于 $\mathbb{R}^{2}-g(A)$ 的无界分支.

因为 $\mathbb{R}^{2}$ 是局部道路连通的, 所以开集 $\mathbb{R}^{2}-g(A)$ 也是局部道路连通的. 从而 $\mathbb{R}^{2}-g(A)$ 的分支和道路分支是相同的. 因此我们可以在 $\mathbb{R}^{2}-g(A)$ 中选取一条从 $0$ 到 $p$ 的道路 $\alpha$. 定义同伦 $G\colon A\times I\to\mathbb{R}^{2}-0$ 为
\[
G(x,t)=g(x)-\alpha(t);
\]
这个同伦画在 \cref{ch10:fig-61-2} 中. 同伦 $G$ 是映射 $g$ 和映射 $k$ 间的同伦, 其中 $k$ 定义为 $k(x)=g(x)-p$. 注意到 $G(x,t)\neq 0$, 因为道路 $\alpha$ 与 $g(A)$ 无交.

现在定义同伦 $H\colon A\times I\to\mathbb{R}^{2}-0$ 为
\[
H(x,t)=tg(x)-p.
\]
$H$ 是映射 $k$ 与某常值映射之间的同伦. 注意 $H(x,t)\neq 0$, 这是因为 $tg(x)$ 在球 $B$ 中, 而 $p$ 却并非如此. 于是我们证明了 $g$ 是零伦的.
\end{proof}

\par
\begin{figure}[htbp]
\centering
\includegraphics[width=0.7\textwidth]{ch10/fig-61-2.png}
\caption{零伦引理的同伦 $G$.}\label{ch10:fig-61-2}
\end{figure}

现在我们来证明 Jordan 分割定理. 设 $X$ 是连通空间, $A\subset X$. 我们称 $A$ 分割 $X$ ($A$ separates $X$), 如果 $X-A$ 不是连通的. 我们称 $A$ 将 $X$ 分割成 $n$ 个分支 ($A$ separates $X$ into $n$ components), 如果 $X-A$ 有 $n$ 个分支.

一段弧 (arc) $A$ 是同胚于单位区间 $[0,1]$ 的一个空间. $A$ 中的两个点 $p$ 和 $q$ 称为 $A$ 的端点 (end point), 如果 $p,q$ 使得 $A-p$ 和 $A-q$ 都是连通的. $A$ 中的其他点就称为 $A$ 的内点 (interior point).

简单闭曲线 (simple closed curve) 是同胚于单位圆周 $S^{1}$ 的空间.

\begin{theorem}[Jordan separation theorem]\label{ch10:thm-jordan-separation}
设 $C$ 是 $S^{2}$ 中的一条简单闭曲线. 则 $C$ 分割 $S^{2}$.
\end{theorem}

\begin{proof}
因为 $S^{2}-C$ 是局部道路连通的, 所以它的分支和道路分支都是局部道路连通的. 以下我们假设 $S^{2}-C$ 是道路连通的, 再由此推出矛盾.

将 $C$ 表示为两段弧 $A_{1}$ 与 $A_{2}$ 之并, 这两段弧仅相交于它们的端点 $a$ 和 $b$. 用 $X$ 表示空间 $S^{2}-a-b$, $U$ 表示 $X$ 中的开集 $S^{2}-A_{1}$, $V$ 表示开集 $S^{2}-A_{2}$. 这时 $X=U\cup V$, 并且
\[
U\cap V=S^{2}-(A_{1}\cup A_{2})=S^{2}-C,
\]
我们已经假设 $U\cap V$ 是道路连通的, 因此满足 Seifert-van Kampen 定理所要求的条件.

设 $x_{0}$ 为 $U\cap V$ 的一个点. 我们来证明内射
\[
i\colon (U,x_{0})\to(X,x_{0})\quad\text{和}\quad j\colon (V,x_{0})\to(X,x_{0})
\]
诱导出所涉及基本群之间的平凡同态. 从而根据 Seifert-van Kampen 定理可见群 $\pi_{1}(X,x_{0})$ 是平凡的. 但 $X=S^{2}-a-b$ 同胚于穿孔平面 $\mathbb{R}^{2}-0$, 因此它的基本群不是平凡的.

我们来证明 $i_{*}$ 是平凡同态. 给定一个基点在 $x_{0}$ 处的回路 $f\colon I\to U$. 为了证明 $i_{*}([f])$ 是平凡的, 令 $p\colon I\to S^{1}$ 是生成 $\pi_{1}(S^{1},b_{0})$ 的标准回路. 映射 $f\colon I\to U$ 诱导连续映射 $h\colon S^{1}\to U$ 使得 $h\circ p=f$. 参见 \cref{ch10:fig-61-3}.

考虑映射 $i\circ h\colon S^{1}\to S^{2}-a-b$. 根据假设, 集合 $i(h(S^{1}))=h(S^{1})$ 与包含点 $a$ 和 $b$ 的连通集 $A_{1}$ 无交. 因此, $a$ 和 $b$ 属于 $S^{2}-i(h(S^{1}))$ 的同一个分支. 根据前述引理, 映射 $i\circ h$ 是零伦的, 因此 $(i\circ h)_{*}$ 是基本群的平凡同态. 但是
\[
(i\circ h)_{*}([p])=[i\circ h\circ p]=[i\circ f]=i_{*}([f]).
\]
因此, $i_{*}([f])$ 是平凡的. 证明完成.
\end{proof}

\par
\begin{figure}[htbp]
\centering
\includegraphics[width=0.7\textwidth]{ch10/fig-61-3.png}
\caption{映射 $f$ 诱导 $h\colon S^{1}\to U$.}\label{ch10:fig-61-3}
\end{figure}

让我们考察一下前面的证明. 在证明中究竟用到简单闭曲线 $C$ 的哪些性质呢? 事实上, 我们只需要 $C$ 能够表示成两个闭的连通子集 $A_{1}$ 和 $A_{2}$ 的并, 其中 $A_{1}$ 和 $A_{2}$ 交于两点 $a$ 和 $b$. 注意到这一点我们便可以得到以下分割定理的推广形式, 这个广义分割定理以后还要用到.

\begin{theorem}[general separation theorem]\label{ch10:thm-general-separation}
设 $A_{1}$ 和 $A_{2}$ 是 $S^{2}$ 的两个连通的闭子集, 并且只交于 $a$ 和 $b$ 两点. 则集合 $C=A_{1}\cup A_{2}$ 分割 $S^{2}$.
\end{theorem}

\begin{proof}
我们首先需要说明的是 $C$ 不等于 $S^{2}$. 在前面的证明中这是显然的. 在目前的情形下 $C\neq S^{2}$ 成立是因为 $S^{2}-a-b$ 是连通的, 但 $C-a-b$ 却不连通. (集合 $A_{i}-a-b$ 构成了 $C-a-b$ 的一个分割.)

剩下的证明和前面那个定理的证明完全一样.
\end{proof}

\begin{exercise}\label{ch10:ex-jordan-separation-1}
举例说明环面上的简单闭曲线有时能够分割环面, 有时不能够分割环面.
\end{exercise}

\begin{exercise}\label{ch10:ex-jordan-separation-2}
设 $A$ 是 $\mathbb{R}^{2}$ 的子集, 由拓扑学家的正弦曲线和从 $(0,-1)$ 到 $(0,-2)$ 到 $(1,-2)$ 到 $(1,\sin 1)$ 的折线路之并组成, 见 \cref{ch10:fig-61-4}. 我们称 $A$ 是闭的拓扑学家的正弦曲线 (closed topologist's sine curve). 证明: 若 $C$ 是 $S^{2}$ 中的一个同胚于闭的拓扑学家的正弦曲线的子空间, 则 $C$ 分割 $S^{2}$.
\end{exercise}

\par
\begin{figure}[htbp]
\centering
\includegraphics[width=0.5\textwidth]{ch10/fig-61-4.png}
\caption{闭的拓扑学家的正弦曲线.}\label{ch10:fig-61-4}
\end{figure}